\subsection{Recovering ideals from extended traces and canonical limits}
\label{sec:trace-ideals}

The final two cases concern reconstruction and coordinate transport. The trace-cone development recovers closed ideals from extended traces, then identifies the images of two ideals attached to a given trace. The mathematical setting is the cone of extended,
lower-semicontinuous traces, studied as an invariant of $C^*$-algebras by
\citet{ElliottRobertSantiago2011}. The selected support results concern this
ideal correspondence. The setting is nonunital complex $C^*$-algebras with compatible order and star structures, allowing nondensely defined extended traces \citep{OAITraceIdeals}.

\paragraph{Recovering the ideal order algebraically.}
An extended trace on $A$, represented by Lean's \src{ExtendedTrace} type,
is a function on positive elements with values in
$[0,\infty]$, together with proofs of vanishing at zero, additivity, positive homogeneity,
traciality, and lower semicontinuity. For a closed two-sided star ideal $I$,
the code constructs the extended trace
\begin{equation}
  \iota_I(a)=
  \begin{cases}0,&a\in I,\\ \infty,&a\notin I,\end{cases}
  \qquad a\geq 0.
  \label{eq:ideal-weight}
\end{equation}
To establish additivity for positive $a,b$, heredity of closed ideals gives
$a+b\in I$ exactly when both summands belong to $I$, proving additivity.
The equivalence $x^*x\in I\Longleftrightarrow x\in I$ and star invariance
give traciality; closed sublevel sets give lower semicontinuity. The existence theorem in \src{FiniteIdeal.lean} supplies all fields of this trace, and the later definition selects the resulting witness.

The reverse construction starts with the finite-square domain of a trace:
\begin{equation}
  \mathcal N_\tau=\{x\in A:\tau(x^*x)<\infty\},
  \qquad I_\tau=\overline{\mathcal N_\tau}^{\norm{\cdot}}.
  \label{eq:finite-ideal}
\end{equation}
The source proves that $\mathcal N_\tau$ is a nonunital star subalgebra
absorbing multiplication on both sides, then takes its topological
closure. This realizes $I_\tau$ as a closed ideal of the original algebra. If
$\tau+\tau=\tau$, every value is either zero or infinity.
Finite-trace positive elements approximate every positive element of
$I_\tau$; lower semicontinuity extends their zero values to that limit,
yielding $\tau=\iota_{I_\tau}$. Conversely,
$I_{\iota_I}=I$. These inverse laws establish
\src{closedIdealEquivIdempotent}, an equivalence between closed ideals
and additive-idempotent traces.

The ideal order is also encoded by addition:
\begin{equation}
  I\subseteq J\quad\Longleftrightarrow\quad
  \iota_I+\iota_J=\iota_I.
  \label{eq:ideal-addition-order}
\end{equation}
Consequently, any additive bijection $e$ of the two extended-trace cones
already induces an order isomorphism $\Phi$ of their closed ideals, with
$\iota_{\Phi(I)}=e(\iota_I)$. Additivity and bijectivity suffice for this algebraic reconstruction. Convergence and scalar compatibility enter in the next step, which identifies the ideals attached to a particular trace.

\paragraph{Identifying the transported ideals by limits.}
The analytic step uses a cut-regularized notion of convergence. Write $(a-\varepsilon)_+$ for the positive cut defined by
continuous functional calculus. The development proves
\begin{equation}
  a\in I_\tau\quad\Longleftrightarrow\quad
  \forall\varepsilon>0,\quad
  \tau\bigl((a-\varepsilon)_+\bigr)<\infty,
  \qquad a\geq0.
  \label{eq:finite-ideal-cuts}
\end{equation}
Its canonical convergence predicate for a filtered family $(\tau_i)$ and
a candidate limit $\sigma$ requires, for every positive $a$ and every
$\varepsilon>0$,
\begin{equation}
  \limsup_i\tau_i\bigl((a-\varepsilon)_+\bigr)
  \leq\sigma(a)\leq\liminf_i\tau_i(a).
  \label{eq:trace-canonical-convergence}
\end{equation}
The two inequalities have different roles: the left regularizes the
argument by a positive cut; the right uses the original element. Their
combination, together with recovery of a lower-semicontinuous trace from
its cuts, gives uniqueness of a canonical limit along a proper filter.\footnote{The convergence predicate is defined in
\href{https://github.com/openai/math/blob/fd4aeeb2ee4fc729c18d98444fed42fd0529eeeb/lean/OAI/Analysis/TraceCone/Basic.lean\#L65-L69}{\src{TraceCone/Basic.lean}, lines 65--69}; uniqueness is proved in
\href{https://github.com/openai/math/blob/fd4aeeb2ee4fc729c18d98444fed42fd0529eeeb/lean/OAI/Analysis/TraceCone/IdealTransport.lean\#L177-L190}{\src{IdealTransport.lean}, lines 177--190}.}

In this sense the source proves
\begin{equation}
  r\tau\longrightarrow\iota_{I_\tau}\quad(r\downarrow0),
  \qquad
  r\tau\longrightarrow\iota_{Z_\tau}\quad(r\to\infty),
  \label{eq:trace-scalar-limits}
\end{equation}
where $Z_\tau$ is a closed ideal whose positive elements are exactly those
on which $\tau$ vanishes. The first limit depends on the closure and cut
characterization in~\eqref{eq:finite-ideal-cuts}. Whenever
$a\in I_\tau$ but $\tau(a)=\infty$, every positive scalar multiple
still has value $(r\tau)(a)=\infty$, although
$\iota_{I_\tau}(a)=0$. Finite values on all positive cuts make the scalar-zero limit compatible with~\eqref{eq:trace-canonical-convergence}. The scalar-infinity limit is first proved pointwise using the trace that is zero on $\tau$'s zero set and infinity elsewhere.

The ideal weights, cut-regularized topology, and these two scalar limits
are classical parts of the trace-cone analysis of
\citet[Sections~3.1--3.2]{ElliottRobertSantiago2011}. The source makes their
formal dependencies explicit: additive structure recovers the ideal
order, while convergence and scalar compatibility identify the two
ideals attached to an arbitrary trace. This separates the algebraic reconstruction from the analytic recognition of the transported objects. Section~\ref{sec:trace-example}
gives a sequence-space example in which the canonical scalar-zero limit
and the pointwise limit differ.

\paragraph{Extending convergence preservation to arbitrary nets.}
The convergence interface makes a further choice about its index type.
The source predicate \src{PreservesCanonicalConvergence}, applied to $e$,
quantifies over filters on the
trace cone itself. The mathematical use, however, includes scalar-indexed
families $g:J\to\mathrm{ExtendedTrace}(A)$ and filters $l$ on $J$.
The source passes to the image filter $g_*l$. Write
$\operatorname{CanConv}(l,f,\sigma)$ for the canonical convergence
condition in~\eqref{eq:trace-canonical-convergence}, along $l$ for the
family $f$ with limit $\sigma$. For a further map $f$ into an
extended-trace cone, the reindexing lemma states
\[
 \operatorname{CanConv}(g_*l,f,\sigma)
 \quad\Longleftrightarrow\quad
 \operatorname{CanConv}(l,f\circ g,\sigma),
\]
because the defining limsup and liminf commute with this filter map.
Applying the preservation interface to the identity family on the trace
cone, then reindexing back, proves preservation for $g$.
The lemma \src{PreservesCanonicalConvergence.net} implements this route.
A condition stated using filters on one carrier therefore supports scalar-indexed families and nets on other index types. The filter-map identity is the reusable bridge between these formulations.\footnote{See
\href{https://github.com/openai/math/blob/fd4aeeb2ee4fc729c18d98444fed42fd0529eeeb/lean/OAI/Analysis/TraceCone/IdealTransport.lean\#L238-L257}{filter reindexing and the net extension, lines 238--257}.}

Finally, a \src{TraceConeEquiv} packages an additive bijection, preservation
of zero and positive scaling, and canonical convergence in both
directions. For each scalar limit, the endpoint proof maps the family
through $e$, rewrites using scalar compatibility, and invokes uniqueness
of canonical limits. Injectivity of $I\mapsto\iota_I$ then yields
\[
  \Phi(I_\tau)=I_{e(\tau)},\qquad
  \Phi(Z_\tau)=Z_{e(\tau)}.
\]
The displayed endpoint proofs use forward convergence preservation;
the inverse field belongs to the larger equivalence interface. The construction thus moves from ideal weights through algebraic reconstruction to analytic recognition by limits. Each stage supplies a particular piece of the identifying equation for the transported ideal. This is a specific use of the
abstraction boundaries familiar from large formalizations
\citep{Perfectoid2020,CommelinTopaz2024}.

\subsection{Named graded maps across all seams of a produced diagram}
\label{sec:all-seams}

The character-variety case proves that values produced by a marked
surface construction satisfy every seam equation. The selected theorem,
\src{producedMarkedValues_allSeams}, starts with an existing field-valued
solution of the original diagram. Its hypotheses include properness of
the diagram, a maximal-rank condition on the facet at the chosen seam,
and positive genus of the parent facet. The conclusion refers to the
specific diagram and values constructed from that solution
\citep{OAIAllSeams}. The key issue is preserving named graded data while
changing coordinates and subdividing the surface.

The diagram consists of surface pieces, called facets, joined along
seams; frames specify the graded identifications at those joins. At the
selected seam, the construction inserts a refined band determined by a
simultaneous splitting of the two input frames. It retains the old
boundary circles, appends a mirror boundary at the parent facet, and
introduces a genus-zero disk for each fresh strip; the parent genus drops by one.
The all-seams theorem checks compatibility of the matrix values on
this resulting combinatorial diagram. This roadmap explains the
inherited, mirror, and branch families used below.\footnote{See the
constructor description in
\href{https://github.com/openai/math/blob/fd4aeeb2ee4fc729c18d98444fed42fd0529eeeb/lean/OAI/AlgebraicGeometry/CharacterVarieties/Cutting/MarkedDiagram.lean\#L79-L104}{\src{Cutting/MarkedDiagram.lean}, lines 79--104}.}

\paragraph{The information in a named graded map.}
The seam equation is stronger than equality of flags. For a graded index
set with grade function $a$ and a frame $e$, let
\[
  F_k^e=e\{u:u_i=0\text{ whenever }a(i)\geq k\},
  \qquad q_k^e(v)=\bigl(e^{-1}v\bigr)|_{\{i:a(i)=k\}}.
\]
Here $q_k^e$ retains the named coordinates of grade $k$. The predicate
$\operatorname{SameFramedFlag}(a,e,f)$ requires
\begin{equation}
  F_k^e=F_k^f,\qquad
  q_k^e(v)=q_k^f(v)\quad\text{for all }v\in F_{k+1}^e,
  \label{eq:named-graded-flag}
\end{equation}
for every $k$. Thus the comparison retains the maps representing the
successive graded pieces, rather than only their dimensions or underlying
subspaces.\footnote{See the definition of
\href{https://github.com/openai/math/blob/fd4aeeb2ee4fc729c18d98444fed42fd0529eeeb/lean/OAI/AlgebraicGeometry/CharacterVarieties/Foundation/Flags.lean\#L155-L160}{\src{SameFramedFlag}}.}
In the produced seam equations, the two frames are obtained
from evaluations of the seam's left and right matrix words. Losing these
named maps during a change of basis would weaken the theorem's meaning.
The one-dimensional frame example in Section~\ref{sec:named-maps}
illustrates this distinction without the surface construction.

\paragraph{Keeping one witness across the construction.}
The intermediate construction records the shared coordinate data needed by the seam predicate. An \src{IdentifiedBand} contains a common frame for a
simultaneous splitting, a rank shape, and row and column identifications
of its graded fibers. Its prefix and named-coordinate laws relate those
identifications to the two input frames. The source first proves that
such a record exists using the named band grades and rank equalities.
It then chooses a witness. The opaque declaration \src{cutGeneratorWitness} names this chosen construction. Its existence-based selection supplies coherent, noncomputable frame data for the later definitions.

The same chosen record supplies the shape, row ranks, column ranks, and marked frames of the produced diagram. Those fields also
determine its punctures and values. Independently choosing compatible
dimensions would not provide these shared identifications. The specification therefore links the dimensions to the actual coordinate identifications throughout the construction. Records and algebraic interfaces are
standard techniques \citep{GarillotEtAl2009,Baanen2022}; the case-specific
payload is the simultaneous preservation of the named graded maps
in~\eqref{eq:named-graded-flag}.

\paragraph{From local seams to the produced values.}
The global seam proof is assembled from eight local families: the chosen
old seam, the other inherited seams, the first and internal original
seams, the last and internal mirror seams, and the positive and negative
branch seams. An equivalence parametrizes the grafted diagram; its surjectivity reduces exhaustiveness to sum, endpoint, and Boolean cases. Consequently, the universal quantifier over new seams follows
from a structural coverage theorem. In the inherited family, coordinate
equivalences identify the new frames with the old named seam frame;
grade-preservation lemmas make these equivalences suitable for the full
predicate. The original solution then supplies the seam equation. The local proof therefore preserves the coordinate maps appearing in the seam equation.

One local producer shows what this coordinate preservation requires.
In \src{AtomicFrames}, ports enumerate finite sets of common labeled
atoms. For a splitting vertex, two routes through the parent and child
enumerations yield the same atom. Injectivity of the parent enumeration
turns that atom identity into equality of the coordinate equivalences.
The source transfers this equality to the comparison matrices, whose
identical frames satisfy \src{SameFramedFlag} by reflexivity.
Merging vertices reuse the splitting construction with reversed local
data. Thus the compatibility starts from a common labeled object and
reaches the named coordinate maps.\footnote{See
\href{https://github.com/openai/math/blob/fd4aeeb2ee4fc729c18d98444fed42fd0529eeeb/lean/OAI/AlgebraicGeometry/CharacterVarieties/Cutting/AtomicFrames.lean\#L179-L227}{the splitting, merging, and frame comparison proofs}.}
For an inherited seam, the later proof assembles six concrete inputs:
grade preservation, the two frame identities, parent-word transport,
child-word transport, and the original seam equation. These inputs expose the coordinate-level information still needed after rank preservation.\footnote{See
\href{https://github.com/openai/math/blob/fd4aeeb2ee4fc729c18d98444fed42fd0529eeeb/lean/OAI/AlgebraicGeometry/CharacterVarieties/Frames/InheritedSeams.lean\#L117-L155}{the inherited-seam proof}.}

There remains a second transport at the endpoint. The general
\src{markedCut_allSeams} theorem proves the eight-family result for the
assembled marked-cut data. The produced values use matrices rebased along
the equality of the selected seam rank and its parent-facet rank.
\src{rebaseUnit} performs this equality-induced change of matrix size.
The endpoint proves that the chosen and rebased generator matrices agree,
uses the mirror-rebasing theorem, and transports the resulting equality
through the side-value constructor and then \src{valuesFromPorts}.
Finally, \src{framedSeamEquations_of_values_eq} transfers the general
seam predicate to the exact produced values. This last transport identifies the values in the theorem with the values returned by the constructor.

The boundary assumptions also distinguish two nearby statements. The
all-seams endpoint has boundary indices $b_0,i_0$, but requires no proof
that this occurrence is the chosen seam. The later
\src{producedMarkedSolution}, which assembles surface, seam, and vertex
equations into a complete solution, additionally assumes
\[
  D.\mathrm{boundarySide}
  \langle D.\mathrm{bandCutParent}(q),b_0,i_0\rangle
  =\langle q,\mathrm{none}\rangle.
\]
Here $D$ is the original diagram, $q$ is the selected seam, and
$\mathrm{none}$ is Lean's constructor for the absent optional index
in a seam-side identifier. The extra hypothesis identifies the chosen
boundary occurrence and is used for the surface component. The distinction identifies the selected theorem's contribution to the larger solution construction. A common witness keeps the coordinates coherent, exhaustive local families cover every seam, and equality transport connects the local proofs to the final produced values.
